\documentclass[11pt]{article}
\usepackage{preamble}
\title{Why a scalloped cold-grounded ice base melts less on average: an exact area-partition below the conductive sublayer}
\author{Abdelrahman Saifelden\\ Alexandria Higher Institute of Engineering and Technology (AIET), Alexandria, Egypt\\ \texttt{Abdulrahman.Saifelden22011411@aiet.edu.eg}}
\date{\today}
\begin{document}
\maketitle

\begin{abstract}


Melt at an ice–water interface is rate-limited by transfer through a thin thermal/haline sublayer — the basis of the standard three-equation melt parameterization (Holland \& Jenkins 1999; McPhee 2008; Kader \& Yaglom 1972). For rough or scalloped ice, the prevailing expectation is that geometry \emph{enhances} this transfer and raises melt (Gilpin et al. 1980; Claudin et al. 2017; Bushuk et al. 2019). Using penalized (Brinkman volume-penalization) DNS/LES of a turbulent cavity between a cold ice ceiling and a bed, we report the \textbf{opposite} in the grounded, cold-walled regime, and give the exact mechanism. Our central result is an \textbf{exact area-partition identity}: the mean basal heat a scalloped wall delivers, relative to a flat wall, decomposes to machine precision as $Nu/Nu_{flat} = C_{reatt} + C_{thick} + C_{rev}$ (reattachment-thin, thickened, and reversed populations). Although lee-side reattachment lifts \emph{local} melt $2–3.3\times $ above its own conduction baseline, the \emph{mean} stays \textbf{below} a flat wall ($Nu/Nu_{flat}<1$, peak $0.922$ at the reference resolution) at every amplitude, because the thickened+reversed lee deficit outweighs the thin-reattachment surplus; a two-moment $(1+CV^{2})$ truncation is falsified ($CV_\delta \sim O(1)$), so the suppression is distribution-dominated, not moment-dominated. The ceiling is \textbf{boundary-condition robust} — relaxing the cold-Dirichlet ice wall to a finite-conductance Robin wall, and the no-slip bed to a Navier slip that raises near-bed speed $\sim 100\times $, both leave the mean unchanged — which pins the limit to the thermal conductive sublayer rather than momentum transport, consistent with the three-equation picture. Two further candidate mechanisms testable in 2-D (intermittent roughness plumes; tidal hydraulic switching) also produce no positive mean signal across three subgrid closures ${none, Smagorinsky, backscatter}$; the double-diffusive case, being intrinsically three-dimensional, is reported here as \textbf{untested in 2-D}, not as a null. We frame the deliverable as a \textbf{regime map}: Type I (grounded, cold-walled) — no mean enhancement; Type II (grounding line, flux boundary) — possible; Type III (open shelf) — yes. Scope is explicit: the candidate sweeps are 2-D (3-D only for the capstone and the counter-gradient flux probe) with Brinkman penalization. A three-seed wall-normal convergence and penalization ($\eta $) study shows the suppression \emph{sign} is robust — $Nu/Nu_{flat}<1$ at every resolution ($ny=48–192$), penalization strength (4\ensuremath{\times} $\eta $), and forcing seed — while its \emph{magnitude} is not grid-converged (it varies non-monotonically across resolution, $0.74–0.99$, with per-grid seed scatter $\pm 0.005$). We therefore claim only the resolution-robust results: the \emph{sign} of the ceiling and the \emph{exact} area-partition identity, which holds to machine precision at every resolution. In this regime, scallop geometry suppresses rather than enhances mean basal melt.


\end{abstract}


\section*{1. Introduction}


Basal melt and the effective pressure it modulates set ice-stream sliding and grounding-line stability (Schoof 2010). Melt at an ice–water interface is rate-limited by molecular transfer through a thin thermal/haline sublayer — the basis of the standard three-equation parameterization (Mellor et al. 1986; Holland \& Jenkins 1999; McPhee 2008; Kader \& Yaglom 1972). For rough or scalloped ice the prevailing expectation, from laboratory and theory, is that geometry \emph{enhances} this transfer and raises melt (Gilpin et al. 1980; Claudin et al. 2017; Bushuk et al. 2019). Operational subglacial models, by contrast, carry no flow-enhanced-melt term in the grounded regime (Röthlisberger 1972; Werder et al. 2013). Whether \emph{resolved cavity turbulence} adds melt above the conduction limit at a cold, grounded ice base is rarely tested with resolved simulation.


We test it with resolved penalized DNS/LES under a pre-stated discipline: each candidate mechanism is stated as a falsifiable expectation \textbf{before} the run, so a null is informative rather than a tuning failure. The result runs against the roughness-enhancement expectation: in the grounded, cold-walled regime, scallop geometry \emph{suppresses} mean melt, and we give the exact mechanism (an area-partition identity) and a regime map for where it applies.


\textbf{Contributions.}

\begin{enumerate}
\item An \textbf{exact area-partition identity} $Nu/Nu_{flat} = C_{reatt} + C_{thick} + C_{rev}$ explaining why a scalloped cold-grounded wall melts \emph{less} on the mean despite $2–3\times $ local lee hot spots — opposite to the roughness-enhancement expectation (§4.3).
\item \textbf{Falsification of the $(1+CV^{2})$ two-moment truncation}: the suppression is distribution-dominated, not moment-dominated (§4.3).
\item \textbf{Boundary-condition robustness} (Robin wall; Navier slip \ensuremath{\times}100) pinning the mean ceiling to the thermal conductive sublayer rather than momentum transport (§5).
\item A \textbf{pre-stated, closure-swept null map} of candidate enhancement mechanisms in the cold-grounded (Type-I) cavity, with double diffusion identified as requiring 3-D and reported as untested here (§3).
\item A \textbf{regime map} (Type I/II/III) scoping the result (§6).
\item The growth-carrying \textbf{local lee-flux law} $R_{max}(a/\lambda )$ — linear, not quadratic, bounded $1.9–4.3$, with separation onset at $a/\lambda \approx 0.11$ (§4.4).
\end{enumerate}


\medskip\hrule\medskip


\section*{2. Methods}


\textbf{2.1 Cavity model.} 2-D/3-D penalized (Brinkman volume-penalization) incompressible DNS/LES of a turbulent cavity between a cold ice ceiling and a bed, with a thermal field and an interfacial melt diagnostic; illustrative real bed transects (Lythe \& Vaughan 2001) supply geometry realism (the area-partition result is geometry-general and does not depend on the specific bed dataset). The interfacial melt obeys a Stefan condition $\rho _{i} L v = q$, $q = -k_{th} \partial \theta /\partial n$ (Stefan 1891; Nye 1953; Crank 1984). We define the flat-wall control, the conduction-relative enhancement $R_{mean} = \langle m_{flow}/m_{cond}\rangle $ (each column normalised by its \textbf{own} flow-off conduction on the same geometry), and the flat-baseline Nusselt ratio $Nu/Nu_{flat} = \langle m_{bump}\rangle /\langle m_{flat}\rangle $ (normalised by the flat wall with flow \textbf{on}). These two ratios have \textbf{different baselines} and are not interchangeable — a distinction that resolves an apparent tension in §4.


\textbf{2.2 Pre-stated go/no-go criteria.} Each candidate's expected signature and its falsifiable threshold are fixed before the run (stated per candidate below). Each is swept across three closures ${none, Smagorinsky, backscatter}$ so a verdict that is closure-independent cannot be a closure artefact.


\textbf{2.3 Type-I regime.} ``Type I'' denotes the grounded, cold-walled, closed cavity: cold-Dirichlet ice ceiling, bounded water layer. The hypothesis under test is that mean melt here is bounded by the \textbf{thermal conductive sublayer} rather than by momentum stagnation. §5 stress-tests this against the boundary conditions themselves; §6 generalises to other regimes.


\textbf{2.4 Provenance.} CPU runs plus Tesla P100 (CuPy 14.0.1) confirmations where noted. GPU-only numbers (counter-gradient $C_{G}$; double-diffusion variance amplification) are tagged as such.


\medskip\hrule\medskip


\section*{3. The null map — three subgrid-driven mechanisms}


\textbf{3.1 Candidate 1 — intermittent plumes from ice-base roughness.} \emph{Pre-stated criterion:} band-limited roughness on the ice base should make interfacial melt \emph{intermittent} (skew $S>0$, excess kurtosis $K>3$, peak/mean $>2$), humped at intermediate $Ri\approx 0.3–0.6$. \emph{Result:} thresholds \textbf{not met} and the statistics are \textbf{flat in $Ri$} (relative spread $<2\%$, closure-independent): pooled $peak/mean\approx 1.8$, $skew\approx -0.75$, $kurt\approx 2.2$. The variability is \textbf{geometric, not plume-driven} — thinner ice columns conduct faster — and switching the flow off entirely ($f_{amp}=0$, pure conduction) reproduces $melt_{mean}$ and \code{peak/\allowbreak{}mean} to \textasciitilde{}1 \%. Mean melt is conduction-limited: $melt_{mean}\approx 2.17\times 10^{-4}$, constant to four significant figures across all $Ri$ and closures. The genuine turbulent signal lives in the interior flux $F_{turb}=\langle v'\theta '\rangle $, not at the interface. (roughness / boundary-layer-plume literature)


\textbf{3.2 Candidate 2 — double-diffusive layering (untested in 2-D).} \emph{Pre-stated criterion:} with $Le=100$ (heat diffuses \textasciitilde{}100\ensuremath{\times} faster than salt), salt fingers should give a \emph{humped} thermal Nusselt $Nu_{T}(R_\rho )$ peaking near $R_\rho \approx 2$. \emph{Result, and why it is not a null:} the 2-D runs show \textbf{no hump} ($Nu_{T}$ falls monotonically $1.44\to 0.73$), but salt fingering and double-diffusive staircases are intrinsically three-dimensional, so a 2-D penalized forced run cannot represent them. We therefore report this mechanism as \textbf{untested here}, not as evidence against double-diffusive enhancement; a 3-D unforced double-diffusive run is required (§7). What the 2-D runs \emph{do} show is a real differential-diffusion signature — $Nu_{S} \gg  Nu_{T}$ ($\sim 30\times $ near $R_\rho =1$) and counter-gradient salt flux at strong stabilisation ($Nu_{S}\approx -26$, $R_\rho =10$), which a positive eddy diffusivity cannot represent — and the Turner flux ratio $\gamma \approx 1/R_\rho $ (to \textasciitilde{}1 \%). (Fig. 1)


\begin{figure}[htbp]
\centering
\includegraphics[width=0.82\linewidth]{figures/48_candidate2_doublediff.png}
\caption{Double-diffusive heat and salt Nusselt numbers ($Nu_{T}$, $Nu_{S}$) versus the density ratio $R_\rho $ in the grounded subglacial cavity. The backscatter-resolving closure resists the conductive-sublayer collapse ($Nu_{T}=0.86$ versus $0.73$ unclosed at $R_\rho =10$); the cold-wall mean-melt ceiling is nonetheless unchanged.}
\end{figure}


\textbf{3.3 Candidate 4 — tidal hydraulic switching.} \emph{Pre-stated criterion:} under a tidal body force a wide cavity should switch between \emph{filled} ($H1\to 1$) and \emph{stratified} ($H1\to 0.2$) states, with time-averaged melt peaking at intermediate $Ri$ (switching frequency $f_{switch}$ maximal there). \emph{Result:} at every $Ri$ the cavity is \textbf{monostable filled} ($H1\approx 0.90\pm 0.01$, $f_{switch}=0$) — no hump. (An earlier ``dead-flow'' null was a numerical under-resolution artefact: $ny=48$ let the Brinkman penalty bleed into the \textasciitilde{}10-row fluid gap; resolving at $ny\ge 64$ restores normal flow, $umax\approx 0.6$, KE up \textasciitilde{}500\ensuremath{\times}, so the null is physical, not a construction artefact.) (subglacial hydraulic-switching / drainage-state lit.)


\begin{quote}
\textbf{Null-map summary.} Neither 2-D-testable subgrid-driven mechanism (roughness plumes, tidal switching) lifts the \emph{mean} interfacial melt above conduction in the Type-I cavity, across all three closures; double diffusion is untested in 2-D (§3.2). The flow does carry real interior structure (counter-gradient flux, $F_{turb}$), but it does not reach the cold wall as net mean melt.
\end{quote}


\medskip\hrule\medskip


\section*{4. The scallop — a positive local signal that is mean-bounded}


Candidate 3 (the scallop / melting instability) is the one mechanism with a positive, flow-driven \emph{interfacial} signal, and the cleanest place to establish the ceiling. GPU-confirmed (Tesla P100).


\textbf{4.1 The pre-stated runaway is sign-reversed — the correct ice physics .} The roughness-growth hypothesis (thin ice \ensuremath{\to} faster melt \ensuremath{\to} more roughness \ensuremath{\to} runaway, $\Lambda >0$) is falsified \emph{with the opposite sign}: the conduction-pinned base \textbf{self-smooths}, $\Lambda \approx -0.265 (<0)$, melt/ice-height anticorrelation $corr(m,y_{ice})\approx -0.95$, closure-independent. This is the expected behaviour of a \emph{melting} (not eroding) surface — the runaway was a false analogy to rock/sediment beds. (The value was briefly corrupted to $+0.24$ by a spectral-filter bug and restored on the GPU path; provenance in \code{glaciers/\allowbreak{}REPORT\_\allowbreak{}CANDIDATE3.\allowbreak{}md}.)


\textbf{4.2 The corrected scallop passes its go/no-go, locally.} A flat melting wall is linearly stable, but a finite-amplitude bump larger than the thermal boundary layer drives lee-side separation \ensuremath{\to} recirculation \ensuremath{\to} reattachment heat-flux enhancement — a stable, self-limiting scallop. The staged probe confirms a coherent, closure-independent enhancement relative to its \emph{own} conduction baseline: $R_{mean}\approx 1.06–1.36$, lee peaks $R_{max}\approx 2–3.3$, scaling monotonically with the mean current ($U^{0.8}–U^{1.5}$). The wavelength sweep shows an interior optimum at $n_{waves}=12$ ($\lambda \approx 2.094$, $\lambda /dx\approx 10.7$): a $+12.5\%$ peak in conduction-relative $R_{mean}$ and a peak $Nu/Nu_{flat}=0.9219$, falling off either side; the conduction-relative enhancement is tight ($R_{mean}=1.1273\pm 0.0044$). (Curl 1966)


\textbf{4.3 ...but the mean is still below a flat wall, and we prove why.} The two ratios look contradictory but are consistent \textbf{by construction}: $R_{mean}>1$ measures bump-melt against its own conduction baseline; $Nu/Nu_{flat}<1$ measures the bump against a \emph{flat wall with flow on}. Flow lifts melt above conduction on the bump, yet the bump still delivers less mean heat than a flat wall — both true at once.


The suppression obeys an \textbf{exact area-partition identity} (no moments, no truncation). Partition the interface relative to the flat-wall conductance $m_{flat}=\langle m_{n},flat\rangle $ into \emph{reattachment} ($m_{n}\ge m_{flat}$, thin sublayer), \emph{thickened} ($0<m_{n}<m_{flat}$), and \emph{reversed} ($m_{n}\le 0$, no physical sublayer). Because the populations tile the interface,


\begin{quote}
$Nu/Nu_{flat} = C_{reatt} + C_{thick} + C_{rev}$, $C_{p} = (1/N) \Sigma _{p} m_{n}/m_{flat}$,
\end{quote}


holding to machine precision ($|C_{sum} - Nu/Nu_{flat}| \lesssim  2\times 10^{-16}$). With $surplus = e_{reatt} > 0$, $deficit = -(e_{thick}+e_{rev}) > 0$ ($e_{p} = C_{p} - f_{p}$), the mechanism is one line: \textbf{$Nu<1$ \ensuremath{\iff} $deficit > surplus$.} Full amplitude sweep (Part-C config):


\begin{center}\small
\begin{tabularx}{\linewidth}{>{\raggedright\arraybackslash}X>{\raggedright\arraybackslash}X>{\raggedright\arraybackslash}X>{\raggedright\arraybackslash}X>{\raggedright\arraybackslash}X>{\raggedright\arraybackslash}X>{\raggedright\arraybackslash}X>{\raggedright\arraybackslash}X>{\raggedright\arraybackslash}X}
\toprule
$a/\lambda $ & \code{Nu/\allowbreak{}Nu\_\allowbreak{}flat} & $f_{reatt}$ & $f_{thick}$ & $f_{rev}$ & surplus & deficit & def>surp & $(\delta _{flat}/\langle \delta _{T}\rangle )(1+CV^{2})$ \\
\midrule
0.05 & 0.895 & 0.477 & 0.523 & 0.000 & 0.112 & 0.217 & \textbf{Yes} & 0.99 \\
0.10 & 0.937 & 0.531 & 0.469 & 0.000 & 0.178 & 0.241 & \textbf{Yes} & 1.64 \\
0.15 & 0.888 & 0.406 & 0.555 & 0.039 & 0.180 & 0.292 & \textbf{Yes} & 2.37 \\
0.20 & 0.897 & 0.281 & 0.672 & 0.047 & 0.211 & 0.315 & \textbf{Yes} & 1.10 \\
0.30 & 0.893 & 0.234 & 0.734 & 0.031 & 0.306 & 0.413 & \textbf{Yes} & 2.09 \\
0.50 & 0.930 & 0.195 & 0.773 & 0.031 & 0.418 & 0.488 & \textbf{Yes} & 0.58 \\
\bottomrule
\end{tabularx}
\end{center}


At every amplitude the reattachment surplus is outweighed by the thickened+reversed deficit, so $Nu<1$; the deficit grows monotonically ($0.22\to 0.49$) as the cavity fills with stagnant lee fluid ($f_{thick}: 0.47\to 0.77$). The thin reattachment patches carry a fat $1/\delta _{T}$ tail (top conductance decile holds $0.17\to 0.48$ of $\langle 1/\delta _{T}\rangle $) but are an area minority. The last column shows the \textbf{falsification} of the two-moment $(1+CV^{2})$ truncation: with $CV_\delta \sim O(1)$ it swings $0.58–2.37$ and wrongly predicts $Nu>1$ at 4 of 6 amplitudes. The suppression is therefore \textbf{distribution-dominated, not moment-dominated} — a sharper physical statement (flow-separation structure of the thermal boundary layer), not a weaker one. (\code{glaciers/\allowbreak{}scallop\_\allowbreak{}g1\_\allowbreak{}populations.\allowbreak{}py}, \code{glaciers/\allowbreak{}scallop\_\allowbreak{}sublayer\_\allowbreak{}probe.\allowbreak{}py}.)


\textbf{4.4 The growth-carrying local flux, quantified.} The \emph{mean} is bounded, but melt \textbf{growth} is carried by the local lee flux, which the area-partition does not fit. From the committed amplitude sweep, the local peak ratio $R_{max}(a/\lambda )$ rises \textbf{roughly linearly} with bump steepness — linear $R^{2}=0.945$, free exponent $0.69$ (sub-quadratic), from $1.9$ ($a/\lambda =0.05$) to $4.3$ ($a/\lambda =0.30$); the $(a/\lambda )^{2}$ closure is rejected for the \emph{local} term too (quadratic $R^{2}=0.40\ll 0.94$), and the origin-respecting steepness law $R_{max}\approx 1+1.7\cdot (2\pi  a/\lambda )$ holds to $R^{2}=0.92$ (seed-robust). The mean conductance, by contrast, is amplitude-\textbf{flat} ($R_{mean}$ saturates \textasciitilde{}1.28; the \code{Nu/\allowbreak{}Nu\_\allowbreak{}flat} deficit is constant, $C\approx 1.11$), and the lee flux \textbf{reverses} ($R_{min}<0$) at $a/\lambda \approx 0.11$ — a separation onset the amplitude-flat mean hides. The honest closure for a unified melt rate is therefore: mean amplitude-flat, growth carried by a bounded, linear-in-steepness local lee flux — not a quadratic ansatz (\code{validation/\allowbreak{}synthetic/\allowbreak{}g6\_\allowbreak{}local\_\allowbreak{}flux\_\allowbreak{}law.\allowbreak{}py}, 6 tests).


\medskip\hrule\medskip


\section*{5. The ceiling is boundary-condition robust}


Every null in §3–§4 used a no-slip bed and a cold-Dirichlet ice wall. Two targeted gates (P100, $n=128$, turbulent) show the ceiling is not an artefact of either:


\begin{itemize}
\item \textbf{Thermal BC (Dirichlet \ensuremath{\to} finite-conductance Robin, $q=h(\theta -\theta _{ice})$):} basal heat absorbed is unchanged flow-on vs flow-off to \textbf{$<0.01\%$}.
\item \textbf{Momentum BC (no-slip \ensuremath{\to} Navier slip):} the near-bed tangential speed grows \textbf{\textasciitilde{}100\ensuremath{\times}} ($u_{tang}$ $0.027 \to  2.68$ from locked to free slip), yet \code{Nu/\allowbreak{}Nu\_\allowbreak{}flat} stays $0.96–0.97$ and \textbf{never exceeds 1} (it slightly falls); ridge melt holds to within 0.4 \%, and the slip parameter $s=1$ reproduces the no-slip trajectory bit-for-bit (including on the GPU CuPy path).
\end{itemize}


Removing the stagnant momentum layer entirely delivers \textbf{no} extra basal heat: the limit is the \textbf{thermal conductive sublayer} set by $\kappa $ and the cold wall, not the momentum stagnation layer. (\code{glaciers/\allowbreak{}subglacial/\allowbreak{}slip\_\allowbreak{}gate.\allowbreak{}py}, \code{glaciers/\allowbreak{}subglacial/\allowbreak{}wall\_\allowbreak{}flux.\allowbreak{}py}; \code{THEORY\_\allowbreak{}CAVITY.\allowbreak{}md} §14.4.)


\textbf{Resolution and penalization robustness of the ceiling.} Because the cavity is Brinkman-penalized, the suppression could in principle be a near-wall penalization-smearing or under-resolution artefact. We test this directly with a wall-normal resolution ladder $ny\in {48,64,96,128,192}$ (sublayer spacing $dy=L_{y}/ny$, $L_{y}=2\pi $) at fixed $\eta =5\times 10^{-5}$, and a penalization sweep $\eta \in {2.5,5,10}\times 10^{-5}$ at $ny=128$, each grid point \textbf{ensembled over three forcing seeds} so that turbulent-realization scatter is separated from genuine grid dependence (single-mode bump $a/\lambda =0.20$, $n_{x}=128$, the same local-normal flux diagnostic as §4.3, with a fixed spinup/measure window held constant across the sweep). Two facts emerge:


\begin{itemize}
\item \textbf{The sign is robust.} $Nu/Nu_{flat}<1$ at \emph{every} seed and \emph{every} grid point (21 configurations): the mean-melt suppression survives a 4\ensuremath{\times} refinement of the wall-normal grid, a 4\ensuremath{\times} change in penalization strength, and the forcing realization. Weakening the penalty shifts the ratio only mildly and monotonically ($\eta : 2.5\to 5\to 10\times 10^{-5}$ gives $Nu/Nu_{flat}: 0.783\to 0.807\to 0.836$, spread $0.053$), in the expected direction — a leakier wall passes slightly more flux — so the ceiling is not a penalization-strength artefact.
\item \textbf{The magnitude is not grid-converged.} Across the resolution ladder \code{Nu/\allowbreak{}Nu\_\allowbreak{}flat} varies non-monotonically — $0.786, 0.740, 0.989, 0.807, 0.869$ for $ny=48\ldots 192$ — with a per-grid seed scatter of only $\pm 0.002–0.007$. Because that scatter is $\sim 30\times $ smaller than the $0.18$ spread across resolution, the variation is a real, reproducible grid dependence, not turbulent noise; $ny=96$ robustly nearly erases the suppression ($0.989$). The precise suppression magnitude is therefore resolution-sensitive and we do \textbf{not} certify it as converged.
\end{itemize}


We accordingly claim only what is resolution-robust: the \emph{sign} of the ceiling ($Nu/Nu_{flat}<1$, melts less on the mean) and the \emph{exact} area-partition identity (§4.3), which holds to machine precision at every resolution. Certifying the magnitude would require a body-fitted (penalization-free) near-wall grid or a uniform near-wall refinement, which we leave open. (\code{glaciers/\allowbreak{}scallop\_\allowbreak{}sublayer\_\allowbreak{}convergence.\allowbreak{}py}, \code{glaciers/\allowbreak{}subglacial/\allowbreak{}sublayer\_\allowbreak{}convergence.\allowbreak{}json}.)


\textbf{Independent corroboration of the flux structure.} A 3-D active-buoyancy LES measures \textbf{counter-gradient} subgrid heat flux ($C_{G} > -1$, growing with stratification): the flux carries the global pressure field's memory, which a down-gradient closure cannot represent (a GPU counter-gradient probe). This is consistent with the null map — the flow has real non-local flux structure, but it does not raise the cold-wall mean melt.


\textbf{Interface coupling number — when the passive-BC limit is exact.} The gates above show the \emph{mean} ceiling is BC-robust empirically; a derived \textbf{interface coupling number} says \emph{why} the cold-Dirichlet wall is the correct BC at surge timescales. Writing the ice flux as the linear response $q_{ice}'(s)=H(s)v'(s)$, the interface velocity is $v'=q_{water}'/(\rho _{i}L+H(s))$, so the passive-BC limit is corrected by $\Lambda (\omega )=|H(i\omega )|/(\rho _{i}L)$: $\Lambda (0)=St\le 0.06$ (DC — ice participates at its latent-heat-limited Stefan weight), $\Lambda \to 0$ at high frequency (ice frozen, passive-adiabatic BC exact), crossing over at the ice clock $\tau _{d}=\kappa /\bar{V}^{2}\sim 10^{3}–10^{4} yr$. Because the surge band (0.02–2 yr) is far faster than $\tau _{d}$, in-band $\Lambda <5\times 10^{-5}\ll St$: the interface is a \textbf{passive BC to <1 \%} for the sliding/melt physics, and ice becomes a participating medium only for millennial forcing (\code{validation/\allowbreak{}synthetic/\allowbreak{}interface\_\allowbreak{}coupling\_\allowbreak{}number.\allowbreak{}py}, 3 tests). (Stefan 1891; Crank 1984)


\medskip\hrule\medskip


\section*{6. The deliverable is a regime map}


The null is \textbf{regime-specific}, and stating the regime is the deliverable:


\begin{center}\small
\begin{tabularx}{\linewidth}{>{\raggedright\arraybackslash}X>{\raggedright\arraybackslash}X>{\raggedright\arraybackslash}X>{\raggedright\arraybackslash}X}
\toprule
Regime & Setting & Boundary conditions & Mean melt enhancement? \\
\midrule
\textbf{Type I} & grounded, cold-walled, closed cavity & cold Dirichlet ice + bounded water & \textbf{No} — conduction-sublayer ceiling ($Nu/Nu_{flat}<1$), BC-robust \\
\textbf{Type II} & grounding line & flux / open boundary & \textbf{Maybe} — the closed-cavity assumption is relaxed \\
\textbf{Type III} & open shelf / ice-shelf cavity & warm inflow, free interface & \textbf{Yes} — different BC class, outside this study \\
\bottomrule
\end{tabularx}
\end{center}


This explains why operational subglacial models carry no flow-enhanced-melt term in the grounded regime, and makes explicit the boundary-condition class (Type II/III) a future enhancement would require. It is a map of \emph{where mechanisms operate and where the boundaries lie}, not a universal melt multiplier.


\medskip\hrule\medskip


\section*{7. Discussion and limitations}


\textbf{What this is.} A pre-stated, closure-swept, reproducible demonstration that mean basal melt in the cold-grounded cavity is conduction-sublayer-limited and boundary-condition robust, with an exact mechanism (area-partition) and a regime map.


\textbf{What this is not.} 2-D for the candidate sweeps (3-D LES for the capstone and the counter-gradient flux probe); penalized (Brinkman) rather than body-fitted geometry; idealised forcing. Double diffusion (§3.2) is untested here: salt fingering is intrinsically 3-D, so the 2-D runs neither confirm nor exclude double-diffusive enhancement. The conductive-sublayer ceiling has been stress-tested for resolution and penalization (§5): its \emph{sign} is robust ($Nu/Nu_{flat}<1$ across $ny=48–192$, 4\ensuremath{\times} $\eta $, and three seeds), but its \emph{magnitude} is not grid-converged (non-monotone $0.74–0.99$), so we report a sign-robust ceiling and an exact area-partition identity, \textbf{not} a certified suppression magnitude. The scallop local enhancement is qualitative; its amplitude coefficients are a companion study. Type II/III are not simulated.


\textbf{Open follow-ons.} A penalization-free (body-fitted) or uniformly wall-refined grid to grid-converge the suppression \emph{magnitude} (the sign is already resolution-robust, §5); 3-D unforced double diffusion; a Type-II grounding-line flux-BC cavity to test the ``maybe''; coupling the $Nu_{S}\gg Nu_{T}$ asymmetry into a melt-rate phenomenology.


\textbf{Connection to rough-wall Rayleigh--B\'enard (a derived consistency relation).} The ceiling $Nu/Nu_{flat}\le 1$ is the \emph{no-enhancement branch} of rough Rayleigh--B\'enard transfer once $Nu/Nu_{flat}$ is decomposed into the wetted-area gain $A=\langle \sqrt{1+y'^{2}}\rangle \approx 1+\pi ^{2}(a/\lambda )^{2}>1$ and a local-flux factor $f_{loc}$. Because the measured $Nu/Nu_{flat}\approx 0.96$--$0.97$ sits \emph{below} the area gain, the implied $f_{loc}=(Nu/Nu_{flat})/A<1$ and the local suppression $1-f_{loc}$ must \emph{rise} with steepness, exactly cancelling the geometric area gain ($A:1.02\to 1.39$, $f_{loc}:0.95\to 0.69$, suppression $0.05\to 0.31$ across $a/\lambda \in [0.05,0.20]$). Enhancement would require the lee/separation flux growth (§G.6) to overcome the conduction-sublayer deficit, $f_{loc}>1/A$ --- a regime crossover we flag as conjectural pending lee-flux data. This places the cold-cavity ceiling on the same area-vs-local-flux competition that governs rough-RB heat transport, and identifies the suppression, not the geometry, as the controlling term.


\medskip\hrule\medskip


\section*{8. Reproduce}


\begin{Verbatim}
pip install -r requirements.txt
python glaciers/run_subglacial.py --bed real --out-dir glaciers/figures # §2 cavity (Figs 31–38)
python glaciers/run_candidate1.py # §3.1 roughness plumes
python ocean/run_candidate2.py # §3.2 double diffusion (Fig 48)
python glaciers/run_candidate4.py # §3.3 hydraulic switching
python glaciers/scallop_battery.py # §4.1–4.2 scallop go/no-go (P100)
python glaciers/scallop_sweep.py # §4.2 wavelength sweep
python glaciers/scallop_g1_populations.py # §4.3 exact area-partition mechanism
python glaciers/scallop_sublayer_probe.py # §4.3 (1+CV²) falsification
python glaciers/scallop_sublayer_convergence.py # §5 resolution + penalization robustness (3 seeds)
# §5 BC robustness: glaciers/subglacial/slip_gate.py, wall_flux.py (P100)
python glaciers/theorem3_cg_gpu_probe.py # §5 counter-gradient C_G (Fig 52, P100)
pytest glaciers/tests ocean/tests -v
\end{Verbatim}

\section*{Data and code availability}


The double-diffusive transport result (Fig. 1), the cavity-field, transfer and counter-gradient figure set, and the area-partition tables regenerate from the scripts in the reproduce section. All data and analysis code are in the public repository at https://github.com/abdh0saifelden2-debug/K.


\section*{References}


Angot, P., Bruneau, C.-H., \& Fabrie, P. (1999). A penalization method to take into account obstacles in incompressible viscous flows. \emph{Numerische Mathematik} 81, 497–520. doi:10.1007/s002110050401


Brinkman, H. C. (1949). A calculation of the viscous force exerted by a flowing fluid on a dense swarm of particles. \emph{Applied Scientific Research} A1, 27–34. doi:10.1007/BF02120313


Bushuk, M., Holland, D. M., Stanton, T. P., Stern, A., \& Gray, C. (2019). Ice scallops: a laboratory investigation of the ice–water interface. \emph{Journal of Fluid Mechanics} 873, 942–976. doi:10.1017/jfm.2019.398


Claudin, P., Durán, O., \& Andreotti, B. (2017). Dissolution instability and roughening transition. \emph{Journal of Fluid Mechanics} 832, R2. doi:10.1017/jfm.2017.711


Crank, J. (1984). \emph{Free and Moving Boundary Problems.} Oxford University Press.


Curl, R. L. (1966). Scallops and flutes. \emph{Transactions of the Cave Research Group of Great Britain} 7, 121–160.


Gilpin, R. R., Hirata, T., \& Cheng, K. C. (1980). Wave formation and heat transfer at an ice-water interface in the presence of a turbulent flow. \emph{Journal of Fluid Mechanics} 99, 619–640. doi:10.1017/S0022112080000791


Holland, D. M., \& Jenkins, A. (1999). Modeling thermodynamic ice–ocean interactions at the base of an ice shelf. \emph{Journal of Physical Oceanography} 29, 1787–1800. doi:10.1175/1520-0485(1999)029<1787:MTIOIA>2.0.CO;2


Kader, B. A., \& Yaglom, A. M. (1972). Heat and mass transfer laws for fully turbulent wall flows. \emph{International Journal of Heat and Mass Transfer} 15, 2329–2351. doi:10.1016/0017-9310(72)90131-7


Lythe, M. B., \& Vaughan, D. G. (2001). BEDMAP: A new ice thickness and subglacial topographic model of Antarctica. \emph{Journal of Geophysical Research: Solid Earth} 106, 11335–11351. doi:10.1029/2000JB900449


McPhee, M. G. (2008). \emph{Air-Ice-Ocean Interaction: Turbulent Ocean Boundary Layer Exchange Processes.} Springer. doi:10.1007/978-0-387-78335-2


Mellor, G. L., McPhee, M. G., \& Steele, M. (1986). Ice-seawater turbulent boundary layer interaction with melting or freezing. \emph{Journal of Physical Oceanography} 16, 1829–1846. doi:10.1175/1520-0485(1986)016<1829:ISTBLI>2.0.CO;2


Nye, J. F. (1953). The flow law of ice from measurements in glacier tunnels, laboratory experiments and the Jungfraufirn borehole experiment. \emph{Proceedings of the Royal Society A} 219, 477–489. doi:10.1098/rspa.1953.0161


Radko, T. (2013). \emph{Double-Diffusive Convection.} Cambridge University Press.


Röthlisberger, H. (1972). Water pressure in intra- and subglacial channels. \emph{Journal of Glaciology} 11, 177–203. doi:10.3189/S0022143000022188


Schoof, C. (2010). Ice-sheet acceleration driven by melt supply variability. \emph{Nature} 468, 803–806. doi:10.1038/nature09618


Stefan, J. (1891). Über die Theorie der Eisbildung, insbesondere über die Eisbildung im Polarmeere. \emph{Annalen der Physik} 278, 269–286.


Turner, J. S. (1973). \emph{Buoyancy Effects in Fluids.} Cambridge University Press.


Werder, M. A., Hewitt, I. J., Schoof, C. G., \& Flowers, G. E. (2013). Modeling channelized and distributed subglacial drainage in two dimensions. \emph{Journal of Geophysical Research: Earth Surface} 118, 2140–2158. doi:10.1002/jgrf.20146


\end{document}
